\begin{helper}
Fix a recorded set \(R\) of tagged red and blue base-edge coordinates.  There
is a terminal candidate set \(U\) and an order \(L\) such that
\[
e\in U
\quad\Longleftrightarrow\quad
e\in\noderef{TwoBiteTerminalCoordinateUniverse}(m)
\hbox{ and }e\notin R.
\tag{1}
\]
The list \(L\) has no repetitions and lists exactly \(U\).  The order is
chosen so that it has the following support-classification property.  If
\[
L=P,e,Q,
\]
then every one-sided coordinate that can support a \(C^+(I)\) same-color
witness for the current coordinate \(e\) lies in \(R\cup P\).  In the red case
\(e=\operatorname{red}(a,b)\), for every red base vertex \(r<a,b\), both
\(\operatorname{red}(r,a)\) and \(\operatorname{red}(r,b)\) lie in
\(R\cup P\).  In the blue case \(e=\operatorname{blue}(a,b)\), the identical
blue-copy assertion holds.
\end{helper}
\begin{proof}
Let \(U=\noderef{TwoBiteTerminalCoordinateUniverse}(m)\setminus R\).  This is
a finite set, and (1) is its definition.  Let \(L\) be the list obtained by
sorting the red elements of \(U\) lexicographically by their ordered endpoint
pair, then sorting the blue elements of \(U\) by the same rule and appending
the blue list after the red list.  Because \(L\) is a sorted listing of a
finite set, it has no repetitions and its entries are exactly \(U\).

Every \(e\in U\) is outside \(R\) by (1).  Also, since
\(\noderef{TwoBiteTerminalCoordinateUniverse}(m)\) contains only increasing
non-loop coordinates, every \(e\in U\) is represented with increasing
endpoint order.

It remains to prove the support-classification property.  Suppose
\(L=P,e,Q\).  Consider the red case \(e=\operatorname{red}(a,b)\).  Since
\(e\) appears in \(L\), it belongs to \(U\), hence \(a<b\).  Let \(r\) satisfy
\(r<a\) and \(r<b\).  Then the coordinates \(\operatorname{red}(r,a)\) and
\(\operatorname{red}(r,b)\) are both increasing non-loop red coordinates, so
both lie in \(\noderef{TwoBiteTerminalCoordinateUniverse}(m)\).

Take either support coordinate, say \(\operatorname{red}(r,a)\).  If it lies
in \(R\), then it is in \(R\cup P\).  If it does not lie in \(R\), then by
(1) it lies in \(U\).  It is lexicographically earlier than
\(\operatorname{red}(a,b)\), because its first endpoint \(r\) is smaller than
the first endpoint \(a\).  Since the red block of \(L\) is sorted
lexicographically and \(L\) has no repetitions, this earlier terminal
coordinate must occur in the prefix \(P\), not at \(e\) and not in the suffix
\(Q\).  Hence it lies in \(R\cup P\).  The same argument applies to
\(\operatorname{red}(r,b)\): if it is unrecorded it lies in \(U\), and it is
lexicographically earlier than \(\operatorname{red}(a,b)\) because its first
endpoint \(r\) is smaller than \(a\).

The blue case is word-for-word the same inside the blue block.  If
\(e=\operatorname{blue}(a,b)\) and \(s<a,b\), then
\(\operatorname{blue}(s,a)\) and \(\operatorname{blue}(s,b)\) are increasing
blue coordinates.  Each is either recorded, and therefore in \(R\), or is an
unrecorded terminal coordinate in \(U\).  In the latter case it is earlier
than \(\operatorname{blue}(a,b)\) in the sorted blue block, so it appears in
the prefix \(P\).  Thus both blue support coordinates lie in \(R\cup P\).
\end{proof}
